\documentclass[11pt]{article}
\usepackage[T1]{fontenc}
\usepackage[utf8]{inputenc}
\usepackage{lmodern}
\usepackage[margin=1in]{geometry}
\usepackage{amsmath,amssymb,amsthm,mathtools}
\usepackage{aliascnt}
\usepackage{microtype}
\usepackage{booktabs,array,longtable}
\usepackage{xcolor}
\usepackage{enumitem}
\usepackage{fancyhdr}
\usepackage{needspace}
\usepackage{hyperref}
\usepackage[nameinlink,noabbrev]{cleveref}
\hypersetup{colorlinks=true,linkcolor=blue!45!black,citecolor=blue!45!black,
 urlcolor=blue!45!black,pdftitle={Cofinality, Factorization, and Bezout Failure for Entire Surcomplex Hahn Functions},
 pdfauthor={OpenAI, research draft prepared for Vladimir Reshetnikov}}
\setlength{\headheight}{14pt}
\pagestyle{fancy}\fancyhf{}
\fancyhead[L]{\small Entire surcomplex Hahn functions}
\fancyhead[R]{\small Cofinality and B\'ezout failure}
\fancyfoot[C]{\thepage}
\setlength{\parskip}{3pt}
\setcounter{tocdepth}{1}
\widowpenalty=10000\clubpenalty=10000
\setlength{\emergencystretch}{2em}
\allowdisplaybreaks[2]
\newtheorem{theorem}{Theorem}[section]
\theoremstyle{plain}
\newaliascnt{lemma}{theorem}
\newtheorem{lemma}[lemma]{Lemma}
\aliascntresetthe{lemma}
\crefname{lemma}{lemma}{lemmas}
\Crefname{lemma}{Lemma}{Lemmas}
\theoremstyle{plain}
\newaliascnt{proposition}{theorem}
\newtheorem{proposition}[proposition]{Proposition}
\aliascntresetthe{proposition}
\crefname{proposition}{proposition}{propositions}
\Crefname{proposition}{Proposition}{Propositions}
\theoremstyle{plain}
\newaliascnt{corollary}{theorem}
\newtheorem{corollary}[corollary]{Corollary}
\aliascntresetthe{corollary}
\crefname{corollary}{corollary}{corollaries}
\Crefname{corollary}{Corollary}{Corollaries}
\theoremstyle{definition}
\newaliascnt{definition}{theorem}
\newtheorem{definition}[definition]{Definition}
\aliascntresetthe{definition}
\crefname{definition}{definition}{definitions}
\Crefname{definition}{Definition}{Definitions}
\theoremstyle{definition}
\newaliascnt{example}{theorem}
\newtheorem{example}[example]{Example}
\aliascntresetthe{example}
\crefname{example}{example}{examples}
\Crefname{example}{Example}{Examples}
\theoremstyle{definition}
\newaliascnt{question}{theorem}
\newtheorem{question}[question]{Question}
\aliascntresetthe{question}
\crefname{question}{question}{questions}
\Crefname{question}{Question}{Questions}
\theoremstyle{remark}
\newaliascnt{remark}{theorem}
\newtheorem{remark}[remark]{Remark}
\aliascntresetthe{remark}
\crefname{remark}{remark}{remarks}
\Crefname{remark}{Remark}{Remarks}
\newcommand{\C}{\mathbb C}\newcommand{\R}{\mathbb R}
\newcommand{\Q}{\mathbb Q}\newcommand{\Z}{\mathbb Z}\newcommand{\N}{\mathbb N}
\newcommand{\No}{\mathbf{No}}
\newcommand{\K}{K_\Gamma}\newcommand{\E}{\mathcal E_\Gamma}
\newcommand{\A}{\mathcal A_\Gamma}\newcommand{\OO}{\mathcal O_\Gamma}
\newcommand{\mm}{\mathfrak m_\Gamma}
\DeclareMathOperator{\supp}{supp}\DeclareMathOperator{\lc}{lc}
\DeclareMathOperator{\ord}{ord}\DeclareMathOperator{\cf}{cf}
\DeclareMathOperator{\Div}{Div}\DeclareMathOperator{\inpol}{in}
\DeclareMathOperator{\Ann}{Ann}
\newcommand{\ggA}{\mathrel{\gg_{\!A}}}
\newcommand{\abs}[1]{\lvert#1\rvert}
\newcommand{\norm}[1]{\lVert#1\rVert}
\newcommand{\res}{\operatorname{res}}
\newcommand{\reposha}{\texttt{4896a2808ce30e01b1c86ae3ba2295a64762d246}}

\title{\textbf{Cofinality, Factorization, and B\'ezout Failure\\
for Entire Surcomplex Hahn Functions}\\[5pt]
\large A classification at arbitrary set-sized valuation rank}
\author{Research draft prepared for Vladimir Reshetnikov\\
\small OpenAI}
\date{21 September 2026}

\begin{document}
\maketitle
\begin{abstract}
Let $\Gamma$ be a divisible ordered abelian group and let
$K_\Gamma=\C((t^\Gamma))$. We study power series whose evaluations are
strongly Hahn-summable at every point of this fixed field. We prove an
exact coefficient criterion: $\sum a_nZ^n$ has this property precisely
when $v(a_n)/n$ tends cofinally to $+\infty$ in $\Gamma$. Consequently,
nonpolynomial entire functions exist exactly when the nonzero group
$\Gamma$ has countable cofinality. A support-controlled preparation
argument yields finite root counts on every valuation ball, a canonical
product for every radially finite divisor, and a complete factorization
by zeros at arbitrary rank.

The principal result is a sharp separation between factorization and
interpolation. The entire-function ring is always a GCD domain. When
$\Gamma$ has countable cofinality but no cofinal cyclic subgroup, we
construct two entire functions with disjoint simple zero sets for which
$Af+Bg=1$ has no entire solution. Thus a two-generated proper ideal can
have no common zero. Conversely, a cofinal cyclic subgroup permits a
rank-one coarsening, arbitrary discrete Hermite interpolation, and the
B\'ezout property. Together with the polynomial case this gives a full
cofinality--rank trichotomy. A sharp change-of-workspace theorem explains
why these nonpolynomial functions need not extend to larger surreal
scales.

The arbitrary-rank classification and explicit obstruction are proposed
original contributions, not a claimed solution of a named published
conjecture. Classical Hahn-field inputs and rank-one antecedents are
identified. Complete mathematical proofs are supplied; they have not
been independently refereed or proof-assistant verified.
\end{abstract}
\newpage
\tableofcontents
\newpage

\section{The question, the answer, and its scope}
\label{sec:intro}

\subsection{Why factorization is not the whole story}
Over an algebraically closed field, finite polynomials with disjoint zero
sets generate the unit ideal. In familiar one-variable analytic theories,
interpolation often extends this fact to entire functions. It is tempting
to infer the same conclusion as soon as an infinite product theorem is
available. The inference is false in the higher-rank setting studied here.
We obtain a ring in which every admissible zero divisor is realized, every
function factors completely according to its zeros, and greatest common
divisors exist, yet two functions without a common zero need not satisfy
a B\'ezout identity.

The obstruction is genuinely infinite and depends on the ordering of
scales. At increasingly large arguments, any one entire function is
restricted by a common lower bound on the valuations of its coefficients.
We place a second sequence of zeros so close to the first that the
reciprocal values required by a B\'ezout identity escape those bounds.
Every finite interpolation problem remains solvable by a polynomial.
It is the simultaneous infinite problem that fails.

\subsection{The analytic category}
Throughout, $\Gamma$ is a \emph{set-sized divisible ordered abelian group},
written additively, and
\[
 \K=\C((t^\Gamma)),\qquad
 v\left(\sum_\eta c_\eta t^\eta\right)=\min\{\eta:c_\eta\ne0\}.
\]
A Hahn series has well-ordered support. We put $v(0)=+\infty$.
The complex coefficient field has the \emph{trivial} valuation: every
nonzero complex coefficient has valuation zero. Thus an infinite ordinary
complex sum at one exponent is not an allowed Hahn coefficient operation.

\begin{definition}\label{def:entire}
The ring of \emph{strongly entire Hahn functions} is
\[
 \E=\left\{F=\sum_{n\ge0}a_nZ^n\in\K[[Z]]:
       (a_nz^n)_{n\ge0}\text{ is strongly Hahn-summable for every }z\in\K
       \right\}.
\]
Here strong summability means that the union of the term supports is
well ordered and each exponent belongs to only finitely many term
supports. The fact that this set is a ring is proved below.
\end{definition}

We use $\N=\{0,1,2,\ldots\}$. The variable $Z$ is an ordinary single indeterminate; its exponents are
nonnegative integers. The Hahn exponents occur in the coefficients.
Nothing is asserted about arbitrary generalized series in $Z$, about
arbitrary local functions for the fine surreal topology, or about the
Berarducci--Mantova derivation. All derivatives in this article hold
$\K$ constant.

\subsection{The classification}
A subset of $\Gamma$ is cofinal if every element of $\Gamma$ is at most
some member of that subset. For $\Gamma\ne0$, its cofinality is at least
$\aleph_0$. We say that $h>0$ is \emph{cofinal-cyclic} when
$\{nh:n\ge1\}$ is cofinal. Such an element exists exactly when the positive
Archimedean classes of $\Gamma$ have a greatest member.

\Needspace{18\baselineskip}
\begin{theorem}[Cofinality--rank trichotomy]\label{thm:main}
Let $\Gamma\ne0$ be a divisible ordered abelian group. Then $\E$ is a GCD
domain, its units are exactly $\K^\times$, and exactly one of the following
cases occurs.
\begin{enumerate}[label=\textnormal{(\roman*)},leftmargin=2em]
\item If $\cf(\Gamma)>\aleph_0$, then $\E=\K[Z]$, a principal ideal domain.
\item If $\Gamma$ has a cofinal-cyclic element, then $\E$ contains
nonpolynomial functions and is a non-Noetherian B\'ezout domain. Arbitrary
finite-order jets on every radially finite set can be interpolated by an
entire function.
\item If $\cf(\Gamma)=\aleph_0$ and $\Gamma$ has no cofinal-cyclic element,
then $\E$ contains nonpolynomial functions and is a non-Noetherian GCD
domain that is not B\'ezout. More strongly, there are $f,g\in\E$ with
disjoint simple zero sets and $(f,g)\ne\E$.
\end{enumerate}
For $\Gamma=0$, the strong-summation convention gives $\mathcal E_0=\C[Z]$.
\end{theorem}

A GCD domain is an integral domain in which every pair has a greatest
common divisor, defined up to a unit. A B\'ezout domain is an integral
domain in which every finitely generated ideal is principal. Case (iii)
therefore separates two properties that coincide for finite polynomial
algebra. This is not a distinction about whether the field itself is
algebraically closed: every $\K$ here is algebraically closed.

The proof is distributed through the article. The coefficient criterion
and cofinality dichotomy are in \cref{sec:growth}; preparation and root
counting are in \cref{sec:prep,sec:zeros}; factorization and GCDs are in
\cref{sec:products}; the negative and positive B\'ezout results are in
\cref{sec:obstruction,sec:positive}. \Cref{sec:classification} assembles
these statements without additional assumptions.

\subsection{Relationship to the repository and the literature}
The repository supplied with the question already contains a rank-one
report over $\C((t^\R))$, including preparation, root counting, and a
partial-theta canonical product \cite{RepoRankOne}. It also contains a
global-divisor report in a different ring,
$\mathcal O(U)((t^\Gamma))$, with ordinary holomorphic coefficient
functions on a common complex domain \cite{RepoDivisors}. Its global
support obstructions are not ignored here; they belong to a different
analytic category. Our argument is not a restatement of its Picard-group
or moving-divisor results.

Classical non-Archimedean factorization is an essential antecedent.
Lazard's work is a foundational reference for one-variable zeros
\cite{Lazard}; Cherry explicitly records the one-variable product
factorization and develops GCD and factorization results in several
variables under a complete real-valued non-Archimedean absolute value
\cite{Cherry}. Hahn-field algebraic closedness is a standard theorem,
used here in the form of Poonen's Corollary~4 \cite{Poonen}.
Support-based operator methods are developed extensively by van der
Hoeven \cite{vdH}. We do not claim any of those facts as new.

\paragraph{Novelty statement.}
The proposed contributions are the arbitrary-rank coefficient and
workspace criteria, the explicit non-B\'ezout construction, and their
combination into \cref{thm:main}. The support-controlled preparation and
product arguments provide a self-contained route to that result; their
rank-one specializations are not new. I did not locate the stated
arbitrary-rank classification or the explicit obstruction in the primary
sources consulted. This is a bounded literature statement, not a
certification of priority. The repository review was targeted at its
catalogue, relevant READMEs, and stated scope, not a line-by-line audit
of every manuscript. No named published open problem is claimed solved.

\section{Hahn supports and the surreal interpretation}
\label{sec:support}

\subsection{Foundational setting}
The proofs can be carried out in ZFC using only the set-sized field
$\K$. To interpret them inside surcomplex numbers, take $\Gamma$ to be
a divisible ordered additive subgroup of $\No$ and send
\begin{equation}\label{eq:embedding}
 \sum_\gamma c_\gamma t^\gamma
 \longmapsto \sum_\gamma c_\gamma\omega^{-\gamma}.
\end{equation}
The real and imaginary parts are Conway normal forms. The standard
normal-form arithmetic makes this an embedding of fields; the real
coefficient subfield lies in $\No$, and its complexification lies in
$\No[i]$ \cite{Conway,Gonshor}. Our explicit example uses a displayed
subgroup of $\No$, so it does not require an unspecified universal
embedding theorem.

The valuation ring and its maximal ideal are
\[
 \OO=\{x\in\K:v(x)\ge0\},\qquad
 \mm=\{x\in\K:v(x)>0\}.
\]
Their residue field is $\C$. An element $x\in\OO$ has a unique expression
$x=c+\varepsilon$, with $c\in\C$ and $\varepsilon\in\mm$. The residue
map is denoted $\res(x)=c$. It is defined on $\OO$, not on all of $\K$.

\subsection{The support facts actually used}
We use the following classical Hahn--Neumann support facts
\cite[Section~3]{Poonen}\cite{vdH}. They are stated explicitly because
every infinite construction below depends on them.

\begin{lemma}[Hahn--Neumann support calculus]\label{lem:neumann}
Let $S,T$ be well-ordered subsets of an ordered abelian group.
\begin{enumerate}[label=\textnormal{(\alph*)},leftmargin=2em]
\item $S\cup T$ and $S+T$ are well ordered. For every $\eta$, there are
only finitely many pairs $(s,t)\in S\times T$ with $s+t=\eta$.
\item If $S\subseteq\Gamma_{>0}$, its additive monoid
$S^*=\{s_1+\cdots+s_r:r\ge0,\ s_i\in S\}$ is well ordered. For each
$\eta$, there are only finitely many finite ordered words in $S$ whose
sum is $\eta$.
\item Consequently, if $u$ is a Hahn series over a commutative
coefficient ring and $v(u)>0$, the geometric series
$\sum_{r\ge0}u^r$ is strongly summable and is the inverse of $1-u$.
\end{enumerate}
\end{lemma}

Part (b) is a support theorem, not an assertion that $nv(u)$ tends
cofinally to infinity. The latter can fail. For example, in
$\Q e_1\oplus\Q e_2$ with $e_2>ne_1$ for every $n$, the Hahn identity
$\sum_{n\ge0}t^{ne_1}=(1-t^{e_1})^{-1}$ is valid but its partial sums
do not converge in the intrinsic valuation topology: no tail has
valuation greater than $e_2$.

\begin{lemma}[Cofinal tails]\label{lem:cofinal-tails}
Suppose $u_n\in\K$ and $v(u_n)\to+\infty$ cofinally in $\Gamma$;
that is, for every $B\in\Gamma$, eventually $v(u_n)>B$. Then
$(u_n)$ is strongly summable. Its sum has valuation at least
$\min_n v(u_n)$ when at least one term is nonzero.
\end{lemma}
\begin{proof}
For a fixed $B$, only finitely many nonzero term supports meet
$(-\infty,B]$. A descending sequence in their union, beginning at
$B$, would therefore lie in a finite union of well-ordered supports,
which is impossible. Likewise, any exponent occurs in only finitely
many term supports. The valuation estimate follows from coefficientwise
addition. A nonempty set of term valuations has a minimum because it
is a subset of the well-ordered support union.
\end{proof}

We use the usual ZFC equivalence between well-ordering of a linear
order and the absence of an infinite strictly descending sequence.
No proper-class support or proper-class summation is involved.

\subsection{Three meanings of an infinite expression}
It is useful to keep separate (1) a strong Hahn sum, (2) convergence in
the intrinsic valuation topology of a fixed $\K$, and (3) convergence
in the topology inherited from all surreal scales. These need not agree.
Our definition starts with (1). The all-point condition will force
(2) for its evaluation series, even though (1) alone does not. We make
no assertion of convergence in (3). The repository's foundational
reports emphasize this distinction \cite{RepoCatalogue,RepoRankOne}.

\section{An exact coefficient criterion and the cofinality dichotomy}
\label{sec:growth}

For $n\ge1$, division of $v(a_n)$ by $n$ is legitimate because $\Gamma$
is divisible and torsion-free. We interpret $(+\infty)/n=+\infty$.

\begin{theorem}[All-point strong summability]\label{thm:criterion}
Let $\Gamma\ne0$ and $F=\sum_{n\ge0}a_nZ^n\in\K[[Z]]$. The following
conditions are equivalent.
\begin{enumerate}[label=\textnormal{(\arabic*)},leftmargin=2em]
\item $F\in\E$.
\item For every $\gamma\in\Gamma$, the family
$(a_nt^{n\gamma})_{n\ge0}$ is strongly summable.
\item For every $h\in\Gamma$, eventually $v(a_n)>nh$.
\item For every $\gamma,B\in\Gamma$, eventually
$v(a_n)+n\gamma>B$.
\end{enumerate}
In particular, strong entire evaluation is equivalent to the cofinal
condition $v(a_n)/n\to+\infty$.
\end{theorem}
\begin{proof}
The implication (1)$\Rightarrow$(2) is immediate by testing $z=t^\gamma$.
The zero series satisfies all four conditions, so assume $F\ne0$.

Suppose (2) holds. Applying it at $\gamma=0$ shows that the union of
the supports of the $a_n$ is well ordered. Thus
\[
 m=\min\{v(a_n):a_n\ne0\}\in\Gamma
\]
exists. Suppose (3) fails. There is $h\in\Gamma$ for which infinitely
many positive integers $n$ satisfy $a_n\ne0$ and $v(a_n)\le nh$.
Choose $q>0$ with $q>|h|$ and $q>|m|$, where $|\cdot|$ is the ordered
group absolute value. By (2), the set of leading exponents
\[
 \{v(a_n)-2nq:a_n\ne0\}
\]
is well ordered. Its minimum is attained at some finite index $k$;
write that minimum as $\mu$. Then
\[
 \mu=v(a_k)-2kq\ge m-2kq>-(2k+1)q.
\]
Choose a violating index $n>2k+1$. It gives
\[
 v(a_n)-2nq\le nh-2nq<-nq<-(2k+1)q<\mu,
\]
a contradiction. This proves (2)$\Rightarrow$(3). Notice that the
argument used a minimum attained at a \emph{finite index}; it did not
assume that the multiples of $q$ are cofinal in $\Gamma$.

Suppose (3) holds and fix $\gamma,B$. Choose $r>0$ and set
$h=-\gamma+|B|+r$. Eventually
\[
 v(a_n)+n\gamma>n(|B|+r)>B
\]
for $n\ge1$. Thus (4) holds.

Finally, fix $z\ne0$ and put $\gamma=v(z)$. Condition (4) implies
$v(a_nz^n)=v(a_n)+n\gamma\to+\infty$. Apply
\cref{lem:cofinal-tails}. At $z=0$ there is only the constant term.
This proves (4)$\Rightarrow$(1).
\end{proof}

\begin{remark}
Condition (3) is stronger than the statement $v(a_n)\to+\infty$.
For example, over $\Gamma=\Q$, the coefficients $a_n=t^n$ tend to zero
in valuation, but $\sum t^nZ^n$ is not entire: at $Z=t^{-1}$ its
terms all have exponent zero. Conversely, $a_n=t^{n^2}$ satisfies (3)
over $\Q$ and $\R$, but not over a group containing an element larger
than every real number.
\end{remark}

\begin{corollary}[Existence of nonpolynomial functions]\label{cor:cofinality}
For $\Gamma\ne0$, the following are equivalent: $\E\ne\K[Z]$;
$\Gamma$ has a countable cofinal subset; and $\cf(\Gamma)=\aleph_0$.
\end{corollary}
\begin{proof}
If infinitely many $a_n$ are nonzero, their valuations divided by their
positive indices give, by \cref{thm:criterion}, a countable cofinal
subset of $\Gamma$. Conversely, choose a positive increasing cofinal
sequence $(q_n)_{n\ge1}$. Then
\[
 1+\sum_{n\ge1}t^{nq_n}Z^n
\]
satisfies the coefficient criterion and is nonpolynomial. A nonzero
ordered group has no finite cofinal subset, because it has no greatest
element.
\end{proof}

\begin{example}[Two different sources of polynomial rigidity]
If $\Gamma=0$, every nonzero coefficient has support $\{0\}$, so strong
summability at $1$ forces a polynomial. If $\Gamma$ has uncountable
cofinality, there are many infinitesimal coefficients, but countably
many coefficient slopes cannot reach all of its scales. The conclusion
is again polynomiality, for a different reason.
\end{example}

\section{A polynomial-coefficient Hahn algebra and preparation}
\label{sec:prep}

\subsection{The algebra that retains support certificates}
Define
\[
 \A=\C[X]((t^\Gamma)).
\]
An element is $A=\sum_{\eta\in S}P_\eta(X)t^\eta$, where $S$ is well
ordered and every $P_\eta$ is a finite polynomial. No common bound on
$\deg P_\eta$ is imposed. Taking the coefficient of each $X^n$ gives
an injective ring map $\A\hookrightarrow\K[[X]]$.
The support facts in \cref{lem:neumann} justify multiplication on both
sides and show that this map respects it.

For $A\ne0$, set $w(A)=\min\{\eta:P_\eta\ne0\}$. The polynomial
$P_{w(A)}$ is its initial polynomial. Since $\C[X]$ is a domain,
$w(AB)=w(A)+w(B)$ and initial polynomials multiply.

\begin{lemma}[Rescaling criterion]\label{lem:rescaling}
A power series $F=\sum a_nZ^n$ belongs to $\E$ if and only if
\[
 F(t^\gamma X)\in\A\quad\hbox{for every }\gamma\in\Gamma.
\]
For one fixed $\gamma$, membership in $\A$ is equivalent to strong
summability of $(a_nt^{n\gamma})_n$.
\end{lemma}
\begin{proof}
The common Hahn support is exactly the union of the supports of the
scaled coefficients. The coefficient at a fixed Hahn exponent is a
polynomial in $X$ exactly when that exponent occurs at only finitely
many indices $n$. The conclusion follows from \cref{thm:criterion}.
\end{proof}

\begin{lemma}[Units and integral substitution]\label{lem:A-units}
An element of $\A$ is a unit if and only if its initial polynomial is
a nonzero constant. Every $A\in\A$ can be evaluated at any $x\in\OO$;
this evaluation is a ring homomorphism. For $x\in\OO$, substitution
$X\mapsto X+x$ is an automorphism of $\A$.
\end{lemma}
\begin{proof}
If the initial polynomial is $c\ne0$ at exponent $\eta$, write
$A=ct^\eta(1+u)$, with $w(u)>0$, and invert $1+u$ using
\cref{lem:neumann}. Conversely, if $AB=1$, the product of the two
initial polynomials is $1$; both are nonzero constants.

For substitution, write $x=c+\varepsilon$, with $c\in\C$ and
$v(\varepsilon)>0$, treating $\varepsilon=0$ separately. For each
$\eta$, expand the finite polynomial $P_\eta(X+c+\varepsilon)$ in
powers of $\varepsilon$. All supports lie in
\[
 S+(\supp\varepsilon)^*.
\]
This set is well ordered. At a fixed exponent there are only finitely
many decompositions from $S+(\supp\varepsilon)^*$, finitely many words
in $\supp\varepsilon$ for each such decomposition, and finitely many
terms from each $P_\eta$. Thus substitution is a legitimate element
of $\A$. Finite polynomial identities and strong regrouping prove
multiplicativity. Its inverse is substitution by $-x$.
Evaluation follows by setting $X=0$ after this substitution.
\end{proof}

\subsection{Preparation without an Archimedean contraction}
\begin{theorem}[Support-controlled preparation]\label{thm:prep}
Let $A\in\A$ satisfy $w(A)=0$, and suppose its initial polynomial
$P\in\C[X]$ is monic of degree $d$. There are unique $W,U$ such that
\[
 A=WU,\qquad
 W=P+R,\quad R\in\mm[X],\quad\deg R<d,\qquad
 U=1+Q,\quad w(Q)>0.
\]
Here $W$ is a monic polynomial of degree $d$ and $U$ is a unit of $\A$.
If $S=\supp_t(A-P)\subseteq\Gamma_{>0}$, both corrections are supported
in $S^*\setminus\{0\}$.
\end{theorem}
\begin{proof}
Let $M=S^*$, a well-ordered monoid. Write $A=P+\sum_{\eta>0}A_\eta t^\eta$,
with $A_\eta=0$ off $S$. Recursively for $\eta\in M\setminus\{0\}$,
define polynomials $R_\eta,Q_\eta$ by ordinary division by $P$:
\begin{align}
 H_\eta&=A_\eta-
   \sum_{\substack{\alpha+\beta=\eta\\\alpha,\beta>0}}
       R_\alpha Q_\beta,\label{eq:prep-rec1}\\
 H_\eta&=P Q_\eta+R_\eta,\qquad\deg R_\eta<d.
 \label{eq:prep-rec2}
\end{align}
The sum is finite by \cref{lem:neumann}, and every index on its right
is smaller than $\eta$. Thus this is a well-founded recursion on a set.
Set $R=\sum R_\eta t^\eta$ and $Q=\sum Q_\eta t^\eta$.
Because the remainders have a uniform degree bound, $R$ is a finite
polynomial over $\mm$; $Q\in\A$. Equations
\eqref{eq:prep-rec1}--\eqref{eq:prep-rec2} give $A=(P+R)(1+Q)$
coefficient by coefficient. The second factor is a unit.

For uniqueness, suppose two pairs differ. The union of their positive
supports is well ordered, so there is a least exponent $\eta$ where
either pair differs. At that exponent, equality of the products gives
\[
 P(Q_\eta-Q'_\eta)+(R_\eta-R'_\eta)=0.
\]
The second summand has degree less than $d$. Uniqueness of polynomial
division forces both differences to vanish, a contradiction. For $d=0$,
the normalization gives $P=1$, $W=1$, and $U=A$.
\end{proof}

This proof is deliberately not a claim that a Banach iteration converges.
It remains valid when no positive element has cofinal integer multiples.
The finite-degree remainder is the mechanism that extracts a finite zero
scheme from an infinite Hahn series.

\section{Entire operations and exact root counting}
\label{sec:zeros}

\subsection{Algebra and Taylor expansions}
\begin{proposition}\label{prop:operations}
$\E$ is a $\K$-subalgebra of $\K[[Z]]$. Evaluation at each $z\in\K$
is a homomorphism. Translation $F(Z)\mapsto F(Z+b)$, differentiation,
and the primitive with zero constant coefficient preserve $\E$.
The translation coefficients are
\begin{equation}\label{eq:taylor}
 [Z^k]F(Z+b)=\sum_{n\ge k}\binom nk a_n b^{n-k}.
\end{equation}
Entire composition $F\circ G$ is also defined and belongs to $\E$.
\end{proposition}
\begin{proof}
The ring assertion follows from \cref{lem:rescaling}: addition and
multiplication preserve every rescaled copy of $\A$. Evaluation can
be performed by choosing $\gamma=v(z)$, writing $z=t^\gamma x$ with
$x\in\OO$, and applying \cref{lem:A-units}.

For translation and a fixed testing scale $\gamma$, expand
$F(b+t^\gamma X)$ into its finite binomial blocks. Put
$\beta=\min\{v(b),\gamma\}$, with $v(0)=+\infty$. Every coefficient
in the $n$th block has valuation at least $v(a_n)+n\beta$, which tends
cofinally to infinity. Each block has finitely many $X$-coefficients.
The support-union and local-finiteness proof of
\cref{lem:cofinal-tails} therefore applies to the blocks with polynomial
coefficients. This gives an element of $\A$ at every scale and proves
\eqref{eq:taylor}. Regrouping the same jointly summable family proves
the evaluation identity and the inverse translation by $-b$.

For differentiation, $v(na_n)=v(a_n)$ for positive $n$, since $n$ is a
nonzero complex coefficient. Shifting the index by one preserves the
criterion in \cref{thm:criterion}. Division of $a_n$ by $n+1$ also
preserves valuation, so termwise integration preserves it as well.

For composition, fix $\gamma$ and let $G_\gamma=G(t^\gamma X)\in\A$.
If $G_\gamma\ne0$, put $\mu=w(G_\gamma)$. Then
$w(a_nG_\gamma^n)=v(a_n)+n\mu\to+\infty$. Hence the sum is strongly
summable in $\A$. Its coefficients define the same formal composite
at every scale, since rescaling commutes with each coefficient's
strong sum. The case $G=0$ is constant. Evaluation and associativity
follow from the same support-controlled regrouping.
\end{proof}

In the differentiation argument one may make the index shift explicit:
for a fixed $h$, choose $q>0$ with $q>|h|$; the bound
$v(a_n)>n(h+q)$ for large $n$ implies the required bounds with $n-1$
or $n+1$ in place of $n$. Thus no hidden Archimedean estimate is used.

\subsection{Initial polynomials on a ball}
For $\gamma\in\Gamma$, define the closed valuation ball
\[
 B_\gamma=\{z\in\K:v(z)\ge\gamma\}.
\]
A larger ball corresponds to a smaller, typically negative,
$\gamma$. For $F\ne0$, put
\begin{align}
 m_\gamma(F)&=\min_{n:a_n\ne0}\{v(a_n)+n\gamma\},\label{eq:weighted}\\
 I_\gamma(F)(X)&=
 \sum_{v(a_n)+n\gamma=m_\gamma(F)}\lc(a_n)X^n.
 \label{eq:initial}
\end{align}
The minimum is attained and the sum is a finite nonzero polynomial by
\cref{thm:criterion}. Write
$d_\gamma(F)=\deg I_\gamma(F)$ and
$e_\gamma(F)=\ord_{X=0}I_\gamma(F)$.

A zero's multiplicity is the order of the translated formal power series
$F(z+Z)\in\K[[Z]]$. The translation identity makes this intrinsic.

\begin{theorem}[Ball, shell, and residue-direction counts]\label{thm:roots}
For $F\in\E\setminus\{0\}$ and $\gamma\in\Gamma$:
\begin{enumerate}[label=\textnormal{(\alph*)},leftmargin=2em]
\item $B_\gamma$ contains exactly $d_\gamma(F)$ zeros of $F$, counted
with multiplicity.
\item The open ball $\{z:v(z)>\gamma\}$ contains exactly
$e_\gamma(F)$ zeros, and the shell $\{z:v(z)=\gamma\}$ contains
$d_\gamma(F)-e_\gamma(F)$ zeros.
\item For each $c\in\C$, the zeros for which
$v(z)\ge\gamma$ and $\res(t^{-\gamma}z)=c$ have total multiplicity
$\ord_{X=c}I_\gamma(F)$.
\end{enumerate}
\end{theorem}
\begin{proof}
Let $c_0$ be the leading coefficient of $I_\gamma(F)$. Apply
\cref{thm:prep} to
\[
 A=\frac{F(t^\gamma X)}{c_0t^{m_\gamma(F)}}.
\]
It gives $A=WU$, with $W\in\OO[X]$ monic of degree $d_\gamma(F)$,
$\res(W)=I_\gamma(F)/c_0$, and $U$ having initial polynomial $1$.
The field $\K$ is algebraically closed by the standard Hahn-field
theorem \cite[Corollary~4]{Poonen}, so $W$ has exactly its degree many
roots in $\K$, counted with multiplicity.

Every root $x$ of $W$ is integral. Indeed, if $v(x)<0$, its leading
term $x^d$ has valuation $dv(x)$, strictly smaller than that of every
lower-degree term, whose coefficient is integral. Such terms cannot
sum to zero. On $\OO$, evaluation of $U$ has residue $1$ and is nonzero.
After translation by an integral root, $U$ has nonzero constant term,
so its formal local power series is a unit. Hence zeros and
multiplicities of $A$ in $\OO$ coincide exactly with those of $W$.

Factor $W=\prod_j(X-x_j)$ and reduce modulo $\mm$. Then
$\res(W)=\prod_j(X-\res(x_j))$. This proves the residue-direction
count. The direction $c=0$ is exactly the open ball, and the nonzero
directions make up the shell. Rescale back to $Z=t^\gamma X$.
\end{proof}

\begin{corollary}[Units, faithfulness, and values]\label{cor:units}
A nonzero entire function without zeros is a nonzero constant.
Distinct elements of $\E$ define distinct functions on $\K$.
Every nonconstant element of $\E$ takes every value in $\K$.
\end{corollary}
\begin{proof}
If a zero-free $F$ has $a_0=0$, it vanishes at zero, a contradiction.
If $a_n\ne0$ for some $n>0$, choose $\delta>0$ and
\[
 \gamma=\frac{v(a_0)-v(a_n)}{n}-\delta.
\]
Then $v(a_n)+n\gamma<v(a_0)$, so the constant coefficient is not an
active term in \eqref{eq:initial}; in particular, its initial polynomial
has positive degree. \Cref{thm:roots} supplies a zero. Therefore $F$
is constant.

For faithfulness, a nonzero $F$ has only finitely many zeros in any
$B_\gamma$, whereas that ball contains infinitely many elements
$t^\gamma c$ with $c\in\C$. Apply this to the difference of two series.
Finally, for any $b\in\K$, a nonconstant $F-b$ cannot be zero-free,
so $F$ takes the value $b$.
\end{proof}

\begin{corollary}[Initial-form stability]\label{cor:stability}
If $F,G\in\E$ satisfy
$w((F-G)(t^\gamma X))>m_\gamma(F)$, then
$m_\gamma(G)=m_\gamma(F)$ and $I_\gamma(G)=I_\gamma(F)$.
All three counts in \cref{thm:roots} are therefore unchanged.
\end{corollary}
\begin{proof}
A perturbation supported strictly above the leading exponent changes
no coefficient at that exponent. Apply \cref{thm:roots}.
\end{proof}

This is a valuation statement, not a theorem about Euclidean circles
in the fine surcomplex topology. Its purpose here is to retain exact
multiplicities in the factorization argument.

\section{Radial divisors and canonical products at arbitrary rank}
\label{sec:products}

\subsection{Which zero sets are admissible?}
\begin{definition}
An effective divisor on $\K$ is a function $D:\K\to\N$.
It is \emph{radially finite} if
\[
 \sum_{z\in B_\gamma}D(z)<\infty\qquad(\gamma\in\Gamma).
\]
The sum means a finite total multiplicity, not a Hahn sum.
Let $\Div^+_{\mathrm{rad}}(\K)$ denote the monoid of these divisors.
\end{definition}

This is a restriction on valuation balls, not on discrete subsets for
the fine surreal topology. The distinction is necessary: fine
discreteness alone would impose no useful control on an infinite
collection in a fixed bounded valuation ball.

\begin{lemma}[Countability and escape]\label{lem:divisor-countable}
An infinite radially finite divisor has countably many points counted
with multiplicity. Its nonzero points can be enumerated as
$(\rho_j)_{j\ge1}$, with repetitions for multiplicities, and
\[
 v(\rho_j)\longrightarrow-\infty
 \quad\hbox{coinitially in }\Gamma.
\]
In particular, an infinite radially finite divisor implies
$\cf(\Gamma)=\aleph_0$.
\end{lemma}
\begin{proof}
Choose countably many distinct occurrences from the infinite divisor.
Their valuations cannot have a common lower bound, for otherwise one
ball would contain infinitely many occurrences. They form a countable
coinitial subset of $\Gamma$. Balls at those scales exhaust $\K$,
and each contains only finitely many divisor occurrences. Their union
therefore contains the whole divisor and is countable. In any
one-to-one enumeration of its occurrences, only finitely many indices
have valuation at least any fixed $\gamma$. This is the asserted escape.
The occurrence at zero, if present, has finite multiplicity and is
removed before enumerating the nonzero points.
\end{proof}

\subsection{Construction and exact factorization}
\begin{theorem}[Canonical product for a radial divisor]\label{thm:product}
Let $D\in\Div^+_{\mathrm{rad}}(\K)$, let $m=D(0)$, and enumerate its
nonzero occurrences as $\rho_j$. The product
\begin{equation}\label{eq:product}
 P_D(Z)=Z^m\prod_j\left(1-\frac Z{\rho_j}\right)
\end{equation}
is a well-defined element of $\E$, independent of the enumeration.
Its divisor is exactly $D$, and the coefficient of $Z^m$ is $1$.
The product is finite when the nonzero divisor is finite.
\end{theorem}
\begin{proof}
Only the infinite case needs a support argument. Put $b_j=\rho_j^{-1}$.
Fix $\gamma\in\Gamma$ and substitute $Z=t^\gamma X$. By
\cref{lem:divisor-countable}, $v(b_j)+\gamma\to+\infty$.
There are only finitely many factors for which $v(b_j)+\gamma\le0$;
set them aside. For the rest, the supports of $b_jt^\gamma$ have a
well-ordered positive union and are locally finite, by
\cref{lem:cofinal-tails}.

Expand the tail product by finite subsets of indices. Every exponent
belongs to the monoid generated by that positive union. For a fixed
exponent, \cref{lem:neumann} gives finitely many words in the support;
local finiteness gives finitely many choices of factor labels carrying
those exponents. Thus only finitely many product terms contribute.
In particular, the coefficient at that exponent is a polynomial in
$X$. The tail product is an element of $\A$ with initial polynomial
$1$, hence is a unit. Multiplying by the finitely many set-aside factors
and $t^{m\gamma}X^m$ gives an element of $\A$.

For each fixed ordinary degree, this construction sums the same
finite-subset coefficient family, multiplied by $t^{k\gamma}$.
Consequently, the resulting elements for different $\gamma$ are
rescalings of one series in $\K[[Z]]$. By \cref{lem:rescaling}, that
series belongs to $\E$. Strong regrouping also proves independence
of the enumeration.

On $B_\gamma$, the set-aside factors give precisely the roots in
that ball, while the remaining unit is nowhere zero on it. Local
translation shows that multiplicities are the displayed factor
multiplicities. Since the balls cover $\K$, the divisor is exactly
$D$. Every factor of the nonzero-root product has constant term $1$,
so the coefficient of $Z^m$ is $1$.
\end{proof}

\begin{theorem}[Factorization by zeros]\label{thm:factorization}
Let $F=\sum a_nZ^n\in\E\setminus\{0\}$ and let
$m=\ord_{Z=0}F$. Its zero divisor is radially finite. If $\rho_j$ are
all its nonzero zeros, repeated with their multiplicities, then
\begin{equation}\label{eq:factorization}
 F(Z)=a_mZ^m\prod_j\left(1-\frac Z{\rho_j}\right).
\end{equation}
Thus the zero divisor determines $F$ up to a nonzero scalar.
\end{theorem}
\begin{proof}
Radial finiteness is \cref{thm:roots}. Let $G=P_{\Div(F)}$ be the
canonical product from \cref{thm:product}. Since $F$ and $G$ have the
same order $m$ at zero, the formal quotient
\[
 Q=\frac{F/Z^m}{G/Z^m}\in\K[[Z]]
\]
exists: $G/Z^m$ has constant coefficient $1$.

We claim that $Q\in\E$. Fix $\gamma$. Preparation of $F(t^\gamma X)$
and $G(t^\gamma X)$ writes them as
\[
 F(t^\gamma X)=c_F W_F U_F,\qquad
 G(t^\gamma X)=c_G W_G U_G,
\]
where $c_F,c_G\in\K^\times$, $U_F,U_G$ are units of $\A$, and
$W_F,W_G$ are monic polynomials all of whose roots lie in $\OO$.
These roots are exactly the normalized roots of $F$ and $G$ in
$B_\gamma$, with multiplicity. Hence $W_F=W_G$. Cancellation in
$\K[[X]]$, including the common factor $X^m$, gives
\[
 Q(t^\gamma X)=\frac{c_F}{c_G}\frac{U_F}{U_G}\in\A.
\]
This holds for every scale, so \cref{lem:rescaling} proves the claim.

The identity $F=GQ$ is now an identity of entire functions.
Orders at any translated center add in the formal power-series ring.
Since $F$ and $G$ have identical zero multiplicities, $Q$ has no zeros.
By \cref{cor:units}, $Q$ is constant. Comparing the coefficient of
$Z^m$ gives $Q=a_m$.
\end{proof}

The quotient step is important. Equal zero sets alone do not justify
dividing two infinite series in the entire category. Preparation on
\emph{each} ball proves that the formal quotient survives every
rescaling. No compactness or transfinite convergence assertion replaces
that argument.

\begin{corollary}[Divisibility, GCDs, and least common multiples]
\label{cor:gcd}
The map $F\mapsto\Div(F)$ induces an isomorphism of commutative monoids
\[
 (\E\setminus\{0\})/\K^\times
 \ \cong\ \Div^+_{\mathrm{rad}}(\K).
\]
For nonzero $F,G\in\E$, $G$ divides $F$ in $\E$ if and only if
$\Div(G)\le\Div(F)$ pointwise. Greatest common divisors and least common
multiples are given by the pointwise minimum and maximum of divisors.
In particular, $\E$ is a GCD domain.
\end{corollary}
\begin{proof}
Surjectivity is \cref{thm:product}; injectivity modulo constants is
\cref{thm:factorization}. Divisors add under products by Taylor
expansion. If the divisor of $G$ is contained in that of $F$, their
difference is radially finite, so its canonical product supplies the
quotient, up to the scalar fixed by \cref{thm:factorization}.
Pointwise minima and maxima of two radial divisors remain radially
finite. They therefore have canonical products with the defining
universal divisibility properties of a GCD and an LCM.
\end{proof}

\begin{corollary}[Finite jets and evaluation ideals]\label{cor:evaluation}
For $x\in\K$ and $r\ge1$, the kernel of the Taylor jet map
\[
 \E\longrightarrow\K[U]/(U^r),\qquad
 F\longmapsto F(x+U)\bmod U^r
\]
is $(Z-x)^r\E$, and the map is surjective. In particular, evaluation
at $x$ has kernel $(Z-x)\E$ and residue field $\K$.
\end{corollary}
\begin{proof}
The vanishing of the jet is exactly a zero of order at least $r$;
apply \cref{cor:gcd}. Every jet is realized by a finite polynomial in
$Z-x$.
\end{proof}

\subsection{What ``factorization'' does not imply}
\begin{proposition}\label{prop:nonnoetherian}
If $\E$ contains a nonpolynomial function, it is not Noetherian and
is not atomic. Its irreducible elements are precisely scalar multiples
of linear polynomials.
\end{proposition}
\begin{proof}
By \cref{thm:factorization}, a nonpolynomial entire function has
infinitely many zeros, with multiplicities. Let $\rho_1,\rho_2,\ldots$
be an escaping sequence of nonzero roots and put
\[
 T_n(Z)=\prod_{j\ge n}(1-Z/\rho_j).
\]
Then $T_n=(1-Z/\rho_n)T_{n+1}$, so
\[
 (T_1)\subsetneq(T_2)\subsetneq(T_3)\subsetneq\cdots.
\]
Strictness follows either by evaluating multiplicities at $\rho_n$
or from \cref{cor:gcd}. Thus the ring is not Noetherian.

A nonunit has a zero. If its divisor has at least two occurrences,
separating one occurrence from the remaining divisor gives a product
of two nonunits. A singleton divisor gives a scalar multiple of one
linear factor and is irreducible. Consequently a function with an
infinite divisor cannot be a \emph{finite} product of irreducibles.
This is the failure of atomicity.
\end{proof}

The canonical product theorem is therefore not a unique-factorization-domain
theorem. Its product is infinite in a precisely controlled Hahn sense;
an algebraic UFD requires finite irreducible factorizations.

\section{The higher-rank obstruction to interpolation and B\'ezout}
\label{sec:obstruction}

\subsection{A sequence of genuinely new Archimedean scales}
For positive $\alpha,\beta\in\Gamma$, write
\[
 \beta\ggA\alpha
 \quad\Longleftrightarrow\quad
 \beta>k\alpha\text{ for every integer }k\ge1.
\]
This is Archimedean domination, not merely a large finite ratio.

\begin{lemma}\label{lem:escaping-classes}
Suppose $\cf(\Gamma)=\aleph_0$ and no element of $\Gamma$ is
cofinal-cyclic. There is a positive cofinal sequence $(\gamma_n)_{n\ge1}$
such that $\gamma_{n+1}\ggA\gamma_n$ for every $n$.
\end{lemma}
\begin{proof}
Choose a positive cofinal sequence $(q_n)$. After selecting
$\gamma_n$, the set $\{k\gamma_n:k\ge1\}$ is not cofinal; hence it
has a strict upper bound $b_n\in\Gamma$. Choose
$\gamma_{n+1}>\max\{b_n,q_{n+1},\gamma_n\}$.
Start with $\gamma_1\ge q_1>0$. The resulting sequence has the
required domination and is cofinal because it dominates the $q_n$.
\end{proof}

\subsection{A universal bound on values at escaping monomials}
The next lemma is the core obstruction. It is short, but it uses both
parts of the definition: a common coefficient-support lower bound and
legitimate evaluation at the large argument.

\begin{lemma}[Archimedean growth barrier]\label{lem:growth-barrier}
Let $(\gamma_n)$ be positive and cofinal, and let
$\lambda_n>0$ satisfy $\lambda_n\ggA\gamma_n$ for every sufficiently
large $n$. For each $H\in\E$, eventually
\begin{equation}\label{eq:growth-barrier}
 v\bigl(H(t^{-\gamma_n})\bigr)>-\lambda_n.
\end{equation}
The value $+\infty$ is allowed when $H$ vanishes at the argument.
\end{lemma}
\begin{proof}
The assertion is immediate for $H=0$. Write $H=\sum_{k\ge0}h_kZ^k$
and let $m=\min_{h_k\ne0}v(h_k)$, which exists by strong summability
at $Z=1$. For large $n$, $\gamma_n>|m|$. Then each nonzero term at
$Z=t^{-\gamma_n}$ satisfies
\[
 v(h_kt^{-k\gamma_n})\ge m-k\gamma_n
      >-(k+1)\gamma_n>-\lambda_n.
\]
The terms' valuations tend cofinally to infinity as $k\to\infty$,
by \cref{thm:criterion}, and their nonempty set has a minimum attained
at a finite $k$. That minimum is strictly greater than $-\lambda_n$.
The valuation of their Hahn sum is at least this minimum, giving
\eqref{eq:growth-barrier}.
\end{proof}

\begin{corollary}[Explicit failure of value interpolation]
\label{cor:no-interpolation}
For the sequence in \cref{lem:escaping-classes}, the prescription
\[
 H(t^{-\gamma_n})=t^{-\gamma_{n+1}}\qquad(n\ge1)
\]
has no solution $H\in\E$. Each finite subcollection of these
prescriptions has a polynomial solution.
\end{corollary}
\begin{proof}
The valuations of the prescribed values are $-\gamma_{n+1}$, in
conflict with \cref{lem:growth-barrier}. The nodes are distinct, so
ordinary finite Lagrange interpolation over the field $\K$ realizes
any finite subcollection.
\end{proof}

This failure is not due to nodes accumulating in a bounded ball.
Their divisor is radially finite and therefore is the divisor of an
entire function by \cref{thm:product}. Realizability of zeros is strictly
weaker than realizability of values.

\subsection{Two functions with no common zero and no B\'ezout identity}
\begin{theorem}[An explicit pointwise B\'ezout failure]\label{thm:no-bezout}
Assume the hypotheses of \cref{lem:escaping-classes} and choose its
sequence $(\gamma_n)$. Define
\begin{align}
 \rho_n&=t^{-\gamma_n},&
 \sigma_n&=t^{-\gamma_n}+t^{\gamma_{n+1}},\label{eq:paired-roots}\\
 f(Z)&=\prod_{n\ge1}(1-Z/\rho_n),&
 g(Z)&=\prod_{n\ge1}(1-Z/\sigma_n).\label{eq:paired-products}
\end{align}
Then $f,g\in\E$, both have constant coefficient $1$, their roots
are simple and disjoint, and
\[
 \gcd(f,g)=1,\qquad (f,g)\ne\E.
\]
In particular, there are no $A,B\in\E$ with $Af+Bg=1$.
\end{theorem}
\begin{proof}
We have $v(\rho_n)=v(\sigma_n)=-\gamma_n$. Thus both root sequences
escape every valuation ball, so \cref{thm:product} defines the two
entire functions. Different indices have different valuations; at the
same index, $\sigma_n-\rho_n=t^{\gamma_{n+1}}\ne0$. Their zero sets
are consequently disjoint and simple. \Cref{cor:gcd} gives their GCD.

We calculate $v(g(\rho_n))$ exactly. For $j<n$, $\rho_n/\sigma_j$
has negative valuation, so
\[
 v(1-\rho_n/\sigma_j)=\gamma_j-\gamma_n.
\]
For $j=n$,
\[
 1-\rho_n/\sigma_n
   =\frac{\sigma_n-\rho_n}{\sigma_n},\qquad
 v(1-\rho_n/\sigma_n)=\gamma_{n+1}+\gamma_n.
\]
For $j>n$, the ratio has positive valuation. The infinite tail product
has residue $1$ and valuation zero, by the positive-support product
argument used in \cref{thm:product}. Therefore
\begin{equation}\label{eq:lambda}
 \lambda_n:=v(g(\rho_n))
   =\gamma_{n+1}+(2-n)\gamma_n+\sum_{j<n}\gamma_j.
\end{equation}
The sum here is finite. Since $0<\gamma_j<\gamma_n$ for $j<n$, the
part of \eqref{eq:lambda} other than $\gamma_{n+1}$ is bounded in
absolute value by a finite multiple of $\gamma_n$. Hence
$\lambda_n\ggA\gamma_n$.

Suppose $Af+Bg=1$. Evaluate at $\rho_n$. The values are well defined
and evaluation is multiplicative, so
\[
 B(\rho_n)=\frac1{g(\rho_n)},\qquad
 v(B(\rho_n))=-\lambda_n.
\]
For all sufficiently large $n$, this contradicts
\cref{lem:growth-barrier}. Thus the ideal is proper.
\end{proof}

The essential cancellation-sensitive step is \eqref{eq:lambda}.
No unproved assertion that nearby zeros force a small value is used:
the contribution of each earlier root is retained, and the infinite
tail is proved to be a unit. The earlier-root terms can lower the
valuation by finite multiples of $\gamma_n$, but cannot offset the
new Archimedean scale $\gamma_{n+1}$.

\begin{corollary}[A finitely generated ideal invisible to point tests]
\label{cor:invisible-ideal}
In the situation of \cref{thm:no-bezout}, there is a proper two-generated
ideal contained in no evaluation maximal ideal $(Z-x)\E$.
Every maximal ideal containing it is therefore not the kernel of
evaluation at any $x\in\K$.
\end{corollary}
\begin{proof}
The ideal $(f,g)$ is proper but $f$ and $g$ do not vanish simultaneously.
By \cref{cor:evaluation}, containment in $(Z-x)\E$ would mean both
vanish at $x$. A maximal ideal containing a proper ideal exists by
Zorn's lemma; none of these maximal ideals is an evaluation ideal.
\end{proof}

This is a failure of the assertion that \emph{every finitely generated
proper ideal has a common zero}. We do not call it a metric corona
counterexample, because no norm lower bound or bounded-function algebra
has been specified. Nor do we claim that non-evaluation maximal ideals
are absent in the rank-one entire ring; the statement concerns a
\emph{two-generated} ideal already lacking a common zero.

\subsection{The missing image of the infinite jet map}
Let $D$ be the simple divisor with nodes $\rho_n$. The evaluation map
\[
 \operatorname{ev}_D:\E\longrightarrow\prod_{n\ge1}\K,
 \qquad H\longmapsto(H(\rho_n))_n
\]
has kernel $f\E$ by \cref{cor:gcd}. It is not surjective by
\cref{cor:no-interpolation}. In particular,
\[
 \E/(f)\hookrightarrow\prod_{n\ge1}\K
\]
is a proper subring. The element $g\bmod f$ maps to an everywhere
nonzero sequence, hence a unit in the product ring, but it is not a
unit in $\E/(f)$. If it had an inverse there, representatives would
give a B\'ezout identity.

Thus the elementary implication ``invertible at each node implies
invertible in the global quotient'' is exactly where unrestricted
infinite Chinese remainder reasoning fails. Every finite projection
of this map is surjective. Neither a finite truncation nor a finite
root calculation can detect the global failure.

\section{The positive case: a cofinal cyclic scale restores interpolation}
\label{sec:positive}

\subsection{A rank-one coarsening with the same entire functions}
It would be incorrect to assert that the preceding obstruction occurs
for every higher-rank group. A group may have several convex levels
and still have a greatest Archimedean class. In that case a single
cofinal scale permits ordinary non-Archimedean analytic methods.

\begin{lemma}[Real-valued coarsening]\label{lem:coarsening}
Suppose $h>0$ and $\{nh:n\ge1\}$ is cofinal in $\Gamma$. There is an
order-preserving homomorphism $\pi:\Gamma\to\R$ with $\pi(h)=1$,
whose kernel is the convex subgroup
\[
 H=\{\gamma\in\Gamma:n|\gamma|<h\text{ for all }n\ge1\}.
\]
The formula
\[
 |a|_\pi=\exp(-\pi(v(a)))\quad(a\ne0),\qquad |0|_\pi=0
\]
defines a nontrivial real-valued non-Archimedean absolute value on
$\K$. Its topology agrees with the intrinsic valuation topology,
and $\K$ is complete for it. Furthermore,
\begin{equation}\label{eq:entire-coarsening}
 \E=\left\{\sum_{n\ge0}a_nZ^n:
           |a_n|_\pi R^n\longrightarrow0\text{ for every real }R>0
     \right\}.
\end{equation}
\end{lemma}
\begin{proof}
Cofinality bounds every $\gamma$ above and below by integer multiples
of $h$. Define
\[
 \pi(\gamma)=\sup\{q\in\Q:qh\le\gamma\}\in\R.
\]
Rational upper and lower approximations show
$\pi(\alpha+\beta)=\pi(\alpha)+\pi(\beta)$ and
$\pi(h)=1$. Monotonicity is immediate. The kernel consists exactly
of the elements infinitesimal relative to $h$, giving the displayed
$H$. In particular, $\pi(\gamma)>n$ implies $\gamma>nh$, and
$\gamma>(n+1)h$ implies $\pi(\gamma)>n$.
The balls specified by integer multiples of $h$ are cofinal among
all valuation neighborhoods of zero. These inequalities show that
the real-valued and ordered-group topologies agree.

For completeness, take a Cauchy sequence for $|\cdot|_\pi$. Choose
a subsequence $(x_{j_n})$ such that
$v(x_{j_{n+1}}-x_{j_n})>nh$. The differences have valuations tending
cofinally to infinity, so their Hahn sum exists by
\cref{lem:cofinal-tails}. Adding it to $x_{j_1}$ gives a limit of the
subsequence; the Cauchy property gives the same limit for the whole
sequence. A metric space is complete when every Cauchy sequence
converges. Algebraic closedness of $\K$ was already supplied by the
Hahn-field theorem.

By \cref{thm:criterion}, the left side of
\eqref{eq:entire-coarsening} is characterized by
$v(a_n)/n\to+\infty$ in $\Gamma$. This is equivalent to
$\pi(v(a_n))/n\to+\infty$ in $\R$: integer multiples of $h$ test the
forward implication, and for the reverse implication choose a real
bound strictly exceeding $\pi(\gamma)$ for a prescribed
$\gamma\in\Gamma$. The usual elementary coefficient test identifies
this last condition with the right side of
\eqref{eq:entire-coarsening}. Explicitly, it implies
$\pi(v(a_n))-n\log R\to+\infty$ for every $R$; conversely, failure
of the ratio condition gives infinitely many ratios at most some
$C$, and then $R=e^{C+1}$ violates the coefficient test.
\end{proof}

This is a \emph{coarsening}, not a valuation-preserving embedding of
$\Gamma$ in $\R$. The kernel $H$ may be nonzero. It is the cofinality
of the topology and the all-radius coefficient test that permit its
use here. Finer residue-direction data are not identified across the
two valuations.

\subsection{A complete Hermite interpolation argument}
The next proof records the damping mechanism that fails in
\cref{sec:obstruction}. It is included instead of importing an
unqualified analytic interpolation theorem.

\begin{theorem}[Arbitrary radial Hermite interpolation]
\label{thm:interpolation}
Assume $\Gamma$ has a cofinal-cyclic element. Let $(x_n)$ be distinct
points of a radially finite set, and assign a positive integer $m_n$
to each point so that the divisor $\sum_nm_n[x_n]$ is radially finite.
For arbitrary jets $J_n\in\K[U]/(U^{m_n})$, there exists $F\in\E$
with
\[
 F(x_n+U)\equiv J_n(U)\pmod{U^{m_n}}
 \qquad\text{for every }n.
\]
\end{theorem}
\begin{proof}
Use $|\cdot|=|\cdot|_\pi$ from \cref{lem:coarsening}. A radially
finite infinite set can be enumerated with $|x_n|\to\infty$.
Indeed, a ball of any fixed real absolute-value radius is contained
in some original valuation ball, and conversely an original ball
has bounded $|\cdot|_\pi$ radius. Handle the finitely many nodes with
$|x_n|\le1$, including zero, first by a finite Hermite polynomial.
Include their factors in every later vanishing product. Reindex the
remaining nodes by $n\ge1$, so $|x_n|>1$ and $|x_n|\to\infty$.

Suppose a polynomial $S_{n-1}$ has the required jets at earlier nodes,
and put
\[
 Q_{n-1}(Z)=\prod_{j<n}(Z-x_j)^{m_j},
\]
including the finitely many initially handled nodes. It is a unit
in the local jet ring at $x_n$. There is therefore a unique jet
\[
 B_n(U)=\frac{J_n(U)-S_{n-1}(x_n+U)}{Q_{n-1}(x_n+U)}
              \pmod{U^{m_n}}.
\]
For a positive integer $N$, define the degree-less-than-$m_n$
polynomial
\begin{equation}\label{eq:interpolation-r}
 r_{n,N}(U)=B_n(U)(1+U/x_n)^{-N}\pmod{U^{m_n}}
\end{equation}
and the polynomial correction
\begin{equation}\label{eq:interpolation-p}
 p_{n,N}(Z)=Q_{n-1}(Z)(Z/x_n)^N r_{n,N}(Z-x_n).
\end{equation}
It vanishes to the required orders at all earlier nodes, and by
\eqref{eq:interpolation-r} its jet at $x_n$ is the required residual.
The parameter $N$ can thus be used only to control its size.

Set $R_n=|x_n|^{1/2}$, so $1<R_n<|x_n|$ and $R_n\to\infty$.
For a polynomial $P=\sum p_kZ^k$, write
$\|P\|_R=\max_k|p_k|R^k$. If $B_n(U)=\sum_{j<m_n}b_jU^j$, the
coefficient of $U^\ell$ in \eqref{eq:interpolation-r} has absolute
value at most
\[
 \max_{j\le\ell}|b_j|\,|x_n|^{j-\ell}.
\]
This bound is independent of $N$: the binomial coefficients
$\binom{-N}{k}$ are integers and have absolute value at most $1$.
Since $R_n<|x_n|$, substitution $U=Z-x_n$ gives
\[
 \|r_{n,N}(Z-x_n)\|_{R_n}
 \le\max_{j<m_n}|b_j|\,|x_n|^j=:C'_n.
\]
Consequently,
\begin{equation}\label{eq:damping}
 \|p_{n,N}\|_{R_n}
 \le \|Q_{n-1}\|_{R_n} C'_n
          \left(\frac{R_n}{|x_n|}\right)^N.
\end{equation}
The right side tends to zero as the ordinary integer $N$ tends to
infinity. Choose $N_n$ so that this norm is at most $2^{-n}$, and put
$S_n=S_{n-1}+p_{n,N_n}$.

For every fixed real $R>0$, eventually $R\le R_n$, so
$\|p_{n,N_n}\|_R\le2^{-n}$. Hence $S_n$ is Cauchy in the Gauss norm
at every radius. The space of series satisfying
$|a_k|R^k\to0$ is complete for that norm: coefficientwise limits
exist by completeness of $\K$, and a uniform approximation by a
finite polynomial bounds the tail of the limit. Thus the limits for
all radii are one common coefficient series $F$, satisfying
\eqref{eq:entire-coarsening}. Therefore $F\in\E$.

Finally, every specified finite jet is preserved. For a fixed node
$x$ and order $m$, choose $R>|x|$. Coefficients of the translated
polynomial up to degree $m-1$ are continuous for $\|\cdot\|_R$,
as follows by the finite binomial formula and its convergent tail.
All sufficiently late $S_n$ have the prescribed jet, so their limit
does too. This proves simultaneous interpolation.
\end{proof}

\begin{theorem}[B\'ezout in the cofinal-cyclic case]\label{thm:bezout}
If $\Gamma$ has a cofinal-cyclic element, then $\E$ is a B\'ezout
domain. In particular, if $f,g\in\E$ have no common zero, there are
$A,B\in\E$ with $Af+Bg=1$.
\end{theorem}
\begin{proof}
First suppose $f,g$ are nonzero and have no common zero. At each zero
$x$ of $f$, let $m_x$ be its multiplicity. Since $g(x)\ne0$, the
formal inverse of $g(x+U)$ exists modulo $U^{m_x}$. By
\cref{thm:interpolation}, choose $B\in\E$ whose jet is that inverse
at every such $x$. The function $1-Bg$ has a zero of order at least
$m_x$ at each zero of $f$. By \cref{cor:gcd}, it is divisible by $f$
in $\E$. Write $1-Bg=Af$.

For arbitrary nonzero $f,g$, divide both by their GCD $d$, whose
existence is \cref{cor:gcd}. The quotients have no common zero, so
their ideal is the unit ideal. Multiplying the resulting identity
by $d$ shows $(f,g)=(d)$. Cases involving zero are immediate.
Induct on the finite number of generators to obtain the B\'ezout
property.
\end{proof}

For such a divisor $D$, the full jet map is therefore surjective:
\[
 \E\big/\bigl(P_D\bigr)
 \ \cong\ \prod_{x\in\supp D}\K[U]/(U^{D(x)}).
\]
The kernel is again given by divisibility, and surjectivity is
\cref{thm:interpolation}. This exact unrestricted Chinese remainder
statement sharply contrasts with the proper image in
\cref{sec:obstruction}.

\section{Classification and concrete surreal examples}
\label{sec:classification}

\subsection{Assembly of the trichotomy}
\begin{proof}[Proof of \cref{thm:main}]
The ring is an integral domain because it is a subring of $\K[[Z]]$.
Its units are exactly the nonzero constants by \cref{cor:units}:
a unit is zero-free, and nonzero constants are invertible. GCDs
exist by \cref{cor:gcd}.

If there is no countable cofinal subset, \cref{cor:cofinality} gives
$\E=\K[Z]$. If there is a cofinal-cyclic element, its positive integer
multiples give a countable cofinal subset, so nonpolynomial functions
exist; \cref{thm:bezout} supplies the B\'ezout property.
If the group has countable cofinality but no cofinal-cyclic element,
\cref{thm:no-bezout} supplies a proper ideal of two coprime functions.
Such an ideal cannot be principal, since any principal generator
would be a common divisor and hence a unit. Non-Noetherianity in
both nonpolynomial cases is \cref{prop:nonnoetherian}.

These cases are mutually exclusive and exhaustive. For $\Gamma=0$,
strong summability at $1$ says that only finitely many coefficient
supports may contain $0$, hence only finitely many coefficients
are nonzero.
\end{proof}

\begin{center}
\small
\begin{tabular}{@{}>{\raggedright\arraybackslash}p{0.34\textwidth}>{\raggedright\arraybackslash}p{0.23\textwidth}>{\raggedright\arraybackslash}p{0.34\textwidth}@{}}
\toprule
Value-group condition & Entire series & Ideal structure\\
\midrule
No countable cofinal subset & Only polynomials & PID\\[3pt]
Cofinal positive integer multiples of one element & Nonpolynomial
series exist & B\'ezout, non-Noetherian\\[3pt]
Countable cofinality, no cofinal cyclic subgroup & Nonpolynomial
series exist & GCD, not B\'ezout; two-generated proper ideal with no common zero\\
\bottomrule
\end{tabular}
\end{center}

\subsection{An explicit rank-infinite subgroup of the surreal field}
Let
\begin{equation}\label{eq:concrete-group}
 \Gamma_\infty=\bigoplus_{j\ge1}\Q e_j,
 \qquad e_j=\omega^j\in\No,
\end{equation}
with the order inherited from $\No$. Every element has finite
support, and the largest index with a nonzero coefficient determines
its sign. Thus $e_{j+1}\ggA e_j$, and $(e_j)$ is cofinal. No element
is cofinal-cyclic: the largest index in its finite support can always
be exceeded.

Using $t^\gamma=\omega^{-\gamma}$, \cref{thm:no-bezout} gives roots
\begin{equation}\label{eq:surreal-roots}
 \rho_n=\omega^{\omega^n},\qquad
 \sigma_n=\omega^{\omega^n}+\omega^{-\omega^{n+1}}.
\end{equation}
The functions
\begin{align*}
 f(Z)&=\prod_{n\ge1}\bigl(1-\omega^{-\omega^n}Z\bigr),\\
 g(Z)&=\prod_{n\ge1}\left(
 1-\frac{Z}{\omega^{\omega^n}+\omega^{-\omega^{n+1}}}\right)
\end{align*}
are entire on the fixed surcomplex Hahn field $K_{\Gamma_\infty}$,
not on all of $\No[i]$. They have real Hahn coefficients and all the
displayed roots are real surreal numbers. The failure of a B\'ezout
identity remains true even when complex Hahn coefficients are allowed
in $A,B$; consequently it also holds when $A,B$ are restricted to
real Hahn coefficients.

The valuation formula becomes
\[
 v(g(\rho_n))=e_{n+1}+(2-n)e_n+\sum_{j<n}e_j.
\]
Its first values are
\begin{align*}
 \lambda_1&=e_2+e_1, & \lambda_2&=e_3+e_1,\\
 \lambda_3&=e_4-e_3+e_2+e_1,
 &\lambda_4&=e_5-2e_4+e_3+e_2+e_1.
\end{align*}
The negative lower-level coefficients do not affect domination:
the coefficient of the highest basis element $e_{n+1}$ is $1$.

For additional coefficient information, the $k$th coefficient of $f$
is
\[
 [Z^k]f=(-1)^k\sum_{1\le j_1<\cdots<j_k}
                         t^{e_{j_1}+\cdots+e_{j_k}},
\]
with leading exponent $e_1+\cdots+e_k$. The minimum comes from the
unique subset $\{1,\ldots,k\}$. Therefore
\[
 v([Z^k]f)=e_1+\cdots+e_k,
\]
and division by the finite integer $k$ still leaves a cofinal sequence
of Archimedean classes. This independently exhibits the coefficient
growth required by \cref{thm:criterion}.

The reciprocal $1/\sigma_j$ has the explicit support-certified form
\[
 \frac1{\sigma_j}
   =t^{e_j}\sum_{r\ge0}(-1)^r t^{r(e_j+e_{j+1})}.
\]
The series is a valid positive-support geometric identity. Its support
need not itself be cofinal in the full rank-infinite group; that is
irrelevant to the legitimacy of this single coefficient. Entirety of
the product is supplied by the cofinal escape of the factor indices
$j$, not by a false claim about $r(e_j+e_{j+1})$.

\subsection{Why finite rank does not give this counterexample}
For $\Gamma=\Q^d$ in lexicographic order, let $h$ lie in its largest
Archimedean class. Integer multiples of $h$ are cofinal, so the entire
ring is B\'ezout by \cref{thm:bezout}, even when $d>1$. The same is
true for $\Q e_1\oplus\Q e_2$ with $e_2\ggA e_1$: take $h=e_2$.
It is \emph{absence of a greatest Archimedean class}, together with
countable cofinality, that forces the negative case.

For a contrasting polynomial example, let
\[
 \Gamma_{\omega_1}=\bigoplus_{\alpha<\omega_1}\Q e_\alpha
\]
with the largest nonzero index dominant. Each countable family of
finite supports is bounded below some countable ordinal. Hence no
countable subset is cofinal, and $\mathcal E_{\Gamma_{\omega_1}}$
consists only of polynomials. This example shows that infinite rank
by itself does not even guarantee nonpolynomial entire functions.

\section{Change of workspace and the limit of all-surreal entirety}
\label{sec:extension}

The adjective ``entire'' must name its field of arguments. A formal
series with coefficients in one Hahn field may be entire there and
fail even strong summability at a single point of a larger field.
The following result gives an exact boundary.

\begin{theorem}[Sharp change-of-workspace criterion]
\label{thm:extension}
Let $\Gamma\subseteq\Delta$ be divisible ordered abelian groups,
and use the natural inclusion $K_\Gamma\subseteq K_\Delta$.
Then
\begin{equation}\label{eq:extension}
 \mathcal E_\Delta\cap K_\Gamma[[Z]]=
 \begin{cases}
   \mathcal E_\Gamma,&\Gamma\text{ is cofinal in }\Delta,\\
   K_\Gamma[Z],&\Gamma\text{ is not cofinal in }\Delta.
 \end{cases}
\end{equation}
In the second case, a single monomial argument in $K_\Delta$ makes
\emph{every} nonpolynomial series with coefficients in $K_\Gamma$
fail strong summability.
\end{theorem}
\begin{proof}
If $\Gamma$ is cofinal in $\Delta$, a sequence of elements of
$\Gamma\cup\{+\infty\}$ tends cofinally to infinity in $\Gamma$
if and only if it does so in $\Delta$. Apply
\cref{thm:criterion} to the coefficient slopes. The zero-group
case is immediate.

If $\Gamma$ is not cofinal, choose $\delta\in\Delta$ strictly larger
than every element of $\Gamma$. For a nonpolynomial
$F=\sum a_nZ^n\in K_\Gamma[[Z]]$, consider evaluation at
$Z=t^{-\delta}$. For any two nonzero coefficients with $m>n$,
\[
 (v(a_m)-m\delta)-(v(a_n)-n\delta)
       =v(a_m)-v(a_n)-(m-n)\delta<0,
\]
since $v(a_m)-v(a_n)\in\Gamma$ and $(m-n)\delta\ge\delta$.
The leading exponents of the nonzero terms therefore form an
infinite strictly descending sequence. Their support union is not
well ordered. No nonpolynomial series over $K_\Gamma$ can be
strongly entire over $K_\Delta$. Polynomials are entire over either
field, proving \eqref{eq:extension}.
\end{proof}

\begin{corollary}[No repair of the counterexample by a Hahn extension]
\label{cor:no-repair}
For the pair $f,g$ in \cref{thm:no-bezout}, let
$\Gamma\subseteq\Delta$ be a divisible ordered-group extension.
If $\Gamma$ is cofinal in $\Delta$, then $f,g\in\mathcal E_\Delta$
and still no $A,B\in\mathcal E_\Delta$ satisfy $Af+Bg=1$.
If $\Gamma$ is not cofinal, neither $f$ nor $g$ is entire on
$K_\Delta$.
\end{corollary}
\begin{proof}
In the cofinal case, $\gamma_n$ remains cofinal in $\Delta$,
and $\gamma_{n+1}\ggA\gamma_n$ is unchanged. The evaluation
formula \eqref{eq:lambda} is unchanged. The proof of
\cref{lem:growth-barrier} now applies to any candidate
$B\in\mathcal E_\Delta$, even when its coefficients lie outside
$K_\Gamma$. Thus the same contradiction holds.
In the noncofinal case, apply \cref{thm:extension} to the two
nonpolynomial series.
\end{proof}

\begin{corollary}[All-surcomplex strong power series are polynomial]
\label{cor:all-no}
Let $F(Z)=\sum_{n\ge0}a_nZ^n$ be an ordinary countable power series
with coefficients in $\No[i]$. If $(a_nz^n)_n$ is strongly
normal-form-summable for every $z\in\No[i]$, only finitely many
$a_n$ are nonzero.
\end{corollary}
\begin{proof}
The union of the normal-form exponent supports of the real and
imaginary parts of the countably many coefficients is a set.
Its divisible additive span $\Gamma\subseteq\No$, with signs
chosen to match $t^\gamma=\omega^{-\gamma}$, is still a set.
All $a_n$ lie in $K_\Gamma$ under \eqref{eq:embedding}.
The set-sized group $\Gamma$ is bounded above by some surreal
$\delta$. Enlarge to its divisible span with $\delta$.
If infinitely many coefficients are nonzero, the descending-leading-
exponent argument of \cref{thm:extension} shows that evaluation at
$z=\omega^\delta$ is not strongly summable.
\end{proof}

The polynomial-rigidity conclusion has antecedents in the repository's
all-scale discussion \cite{RepoCatalogue,RepoRankOne}; it is not
claimed here as an independently new theorem. The additional statement
is the exact cofinal/noncofinal criterion for a specified pair of
set-sized workspaces and the persistence of the B\'ezout obstruction
under every cofinal extension.

For the concrete group \eqref{eq:concrete-group}, one may take
$\delta=\omega^\omega$. It dominates all finite rational combinations
of the $\omega^j$. Hence the point
$Z=\omega^{\omega^\omega}$ already prevents either nonpolynomial
product from being an all-surcomplex strongly entire function.
There is no contradiction: they are entire on the field explicitly
named in their definition.

\section{Consequences, limitations, and further precise questions}
\label{sec:outlook}

\subsection{What has been established}
The results isolate three logically different phenomena.
First, coefficient slopes determine whether an ordinary countable
power series can be evaluated at every scale in a fixed field.
Second, support-controlled preparation and canonical products
completely describe zeros and divisibility within that category.
Third, the order structure above those scales determines whether
arbitrary values and jets can be prescribed and whether coprime
functions generate the unit ideal.

In particular, the factorization monoid alone cannot decide the
B\'ezout property. In both nonpolynomial cases, a function up to a
scalar is exactly a radially finite zero divisor; every effective
subdivisor is again principal. Nevertheless, only one of those
cases admits unrestricted discrete interpolation. The additive
structure of the ring, not merely its multiplicative divisibility
relation, is essential.

The non-B\'ezout example is not a delicate choice of a transcendental
ordinary complex function or of an uncomputable coefficient. Its
nodes, displacements, and factor coefficients are explicit Hahn
expressions over rational data. What cannot be replaced by a finite
calculation is the quantification over all coefficient indices and
all Archimedean scales.

\subsection{Scope that must not be silently enlarged}
All supports and index sets in this article are sets. The field is a
full Hahn field with complex residue coefficients, and $\Gamma$ is
divisible. Algebraic closedness is used to split the prepared finite
polynomial. A nondivisible group or a restricted subfield of Hahn
series may lose the required roots, so the same factorization theorem
cannot simply be copied there.

The ring is one-variable and strong-summation entire. It is not the
ring of ordinary-coefficient holomorphic Hahn families on one common
complex domain, not a ring of arbitrary fine-local surcomplex germs,
and not a several-variable analytic algebra. The arbitrary-rank
preparation theorem is proved in $\C[X]((t^\Gamma))$, not asserted in
a larger power-series coefficient ring. The rank-one coarsening in
\cref{sec:positive} changes residue information, even though it preserves
the topology and the entire-function ring.

There is no construction of a contour integral, a global surcomplex
exponential at unrestricted imaginary arguments, or an oscillator for
a surreal derivation. The result about primitives in
\cref{prop:operations} concerns termwise $Z$-integration with
coefficients fixed. It says nothing about solving differential
equations in the coefficient field with the Berarducci--Mantova
derivation.

\subsection{Questions opened by the classification}
The following are research questions generated by this article, not
claims that a particular publication already lists them as open.

\begin{question}[Exact interpolation image]\label{q:image}
When $\Gamma$ has countable cofinality but no greatest Archimedean
class, characterize the image of the map
\[
 \E\longrightarrow\prod_{x\in\supp D}\K[U]/(U^{D(x)})
\]
for a prescribed radially finite divisor $D$.
\end{question}
The growth barrier gives explicit necessary inequalities and proves
that the image can be proper. It is not proved to be a sufficient
criterion for arbitrary nodes or jets. A satisfactory answer should
encode the compatibility of divided differences across convex
subgroups, rather than only the size of individual values.

\begin{question}[Finitely generated ideals beyond their common divisor]
Can one describe the ideal $(f,g)$ by its common zero divisor together
with an explicit obstruction in an interpolation quotient?
\end{question}
There is already an exact reduction: after dividing the GCD, the
B\'ezout problem is whether the inverse jets of one function along
the other's zero divisor belong to the interpolation image. The
remaining issue is a usable intrinsic description of that image.
The counterexample proves that a zero-divisor description alone is
insufficient.

\begin{question}[Several variables]
Which parts of the cofinality criterion and the B\'ezout obstruction
extend to several ordinary variables, and which local preparation
statements retain explicit Hahn support certificates?
\end{question}
One must separate the cofinal growth of total degree from the geometry
of positive-dimensional zero sets. Cherry's several-variable work is
an important rank-one antecedent \cite{Cherry}; the present proof does
not establish an arbitrary-rank multivariable GCD theorem.

\subsection{A proof-assistant route}
The research statement does not require a formalization of the entire
proper class $\No$. A first formal development can work over a
set-sized divisible linearly ordered abelian group and a Hahn field.
The critical reusable pieces are: the monomial-testing criterion,
the lower-bound and attained-minimum lemmas, preparation by
well-founded support recursion, and the finite root-count bridge.
The explicit obstruction itself then requires only ordered-group
inequalities, finite product valuations, and evaluation homomorphisms.

The supplied repository already exposes Hahn support lemmas and
finite initial-polynomial root calculations in its reported
formalization scope \cite{RepoRoot}. Those modules are plausible
prerequisites, not a certificate that the new proofs have been
formalized. No Lean code or kernel-check claim is included in this
package.

\appendix
\section{Proof-dependency and hypothesis audit}
\label{app:audit}

The main proof can be read along the following dependency chain:
\[
 \begin{gathered}
 \text{Hahn support calculus}\\
 \Downarrow\\[-2pt]
 \text{cofinal coefficient criterion and the algebra }\A\\
 \Downarrow\\[-2pt]
 \text{support preparation and finite root counts}\\
 \Downarrow\\[-2pt]
 \text{canonical products, divisibility, and GCDs}\\
 \Downarrow\\[-2pt]
 \begin{array}{c@{\qquad\qquad}c}
 \text{new Archimedean scales}&\text{one cofinal cyclic scale}\\
 \Downarrow&\Downarrow\\
 \text{growth barrier and no B\'ezout}&\text{Hermite interpolation and B\'ezout}.
 \end{array}
 \end{gathered}
\]
The arrow from preparation to root counting additionally uses
algebraic closedness of the Hahn field. The arrow to interpolation
uses completeness of the real-valued coarsening, proved directly in
\cref{lem:coarsening}. No proof of the negative result depends on the
positive interpolation theorem.

\paragraph{Attained minima.}
The proof of \cref{thm:criterion} uses a minimum attained at an ordinary
finite coefficient index. The proof of \cref{lem:growth-barrier} also
uses an attained minimum; termwise strict inequalities alone would not
justify a strict inequality for an arbitrary limiting process.

\paragraph{Two independent infinities.}
The finite index $k$ in a power series is not a Hahn exponent. Integer
multiples of a positive exponent need not be cofinal. The support
lemma handles geometric-series expansions without that property;
the positive interpolation proof uses it only after constructing a
legitimate real-valued coarsening.

\paragraph{Divisibility is not assumed.}
The quotient $F/G$ is first formed after removing the shared zero at
zero, in $\K[[Z]]$. Only preparation at every scale proves it entire.
This avoids assuming the global divisibility theorem in its own proof.

\paragraph{The counterexample is infinite.}
No finite truncation of the pair can witness non-B\'ezout behavior:
the finite truncations are polynomials with disjoint root sets and
therefore generate the unit ideal. The contradiction requires a
single candidate coefficient lower bound to work at cofinally many
successively dominating scales.

\section{Reproducible finite checks}
\label{app:checks}

The archive includes \texttt{verify\_examples.py}, an exact-arithmetic
Python program using only the standard library. It checks finite
instances of the reverse-lexicographic exponent inequalities, the
valuation formula \eqref{eq:lambda}, leading exponents of finite
canonical products, and a finite-truncation implementation of the
preparation recursion \eqref{eq:prep-rec1}--\eqref{eq:prep-rec2}.
It also checks the local inverse-jet identity that underlies
\eqref{eq:interpolation-r}.

The delivered run passed all \textbf{831 finite checks}; the category counts are
72 paired-root checks, 40 finite dominance samples, 88 product-coefficient
checks, 456 preparation checks, and 175 inverse-jet checks.
The run record is \texttt{verification.json}. Its count is a count of
finite assertions, not a confidence score for the infinite theorem.
All group arithmetic uses rational coordinates with the largest
nonzero coordinate dominant; no floating-point approximation is
used. In the preparation check, the $t$-adic truncation is explicitly
finite, so higher coefficients are discarded only after the claimed
congruence has been specified.

These checks can detect an indexing error, a sign error in the paired
root valuation, or a wrong recurrence coefficient. They cannot prove
cofinality, Hahn support well-ordering for every permitted group,
the infinite product theorem, or nonexistence of a global B\'ezout
solution. Those conclusions rest on the mathematical proofs in the
body of the article.

\section{Source and novelty ledger}
\label{app:sources}

The repository commit used for the final source check was
\begin{center}\small\reposha.\end{center}
The review used the root README, the documentation catalogue, and
the READMEs for the rank-one and global-divisor reports. These four
files were re-read at the displayed commit after the working branch
changed during preparation. The root
README distinguishes implemented finite/Hahn statements from remaining
foundational implementation work. We do not reinterpret missing
implementation as an unsolved mathematical problem.

The primary literature checked included Poonen's Hahn-field
construction and algebraic-closedness corollary, Cherry's complete
real-valued non-Archimedean setting and one-variable factorization
statement, Lazard's bibliographic record as the historical reference,
and van der Hoeven's generalized-series operator framework. The
Conway and Gonshor books, with bibliographic cross-checks against
\cite{ADH}, are cited for the standard surreal normal-form
background, not claimed to have been re-audited in full during this
work. General web searches were also attempted, but several returned
irrelevant results and were not treated as evidence that a theorem
is absent from the literature.

The strongest proposed original statement is \cref{thm:no-bezout}
and its use in \cref{thm:main}. The sharp workspace criterion and
persistence result are \cref{thm:extension,cor:no-repair}.
The rank-one factorization formula is explicitly classical. The
all-surcomplex polynomial corollary is credited to the pre-existing
all-scale theme in the supplied repository. A later discovery of
an antecedent would change the priority discussion, not the content
or hypotheses of the displayed proofs.

\clearpage
\phantomsection\addcontentsline{toc}{section}{References}
\begin{thebibliography}{99}
\bibitem{Conway}
J. H. Conway, \emph{On Numbers and Games}, London Mathematical Society
Monographs, vol.~6, Academic Press, 1976.

\bibitem{Gonshor}
H. Gonshor, \emph{An Introduction to the Theory of Surreal Numbers},
London Mathematical Society Lecture Note Series, vol.~110,
Cambridge University Press, 1986.

\bibitem{Poonen}
B. Poonen, Maximally complete fields, \emph{L'Enseignement
Math\'ematique} (2) \textbf{39} (1993), 87--106.
Author's text: \url{https://math.mit.edu/~poonen/papers/amsval.pdf}.
Section~3 and Corollary~4 are the principal inputs used here.

\bibitem{Lazard}
M. Lazard, Les z\'eros d'une fonction analytique d'une variable sur un
corps valu\'e complet, \emph{Publications Math\'ematiques de l'IH\'ES}
\textbf{14} (1962), 47--75. DOI: \href{https://doi.org/10.1007/BF02684326}{10.1007/BF02684326}.
\url{https://www.numdam.org/item/PMIHES_1962__14__47_0/}.

\bibitem{Cherry}
W. Cherry, Existence of GCD's and factorization in rings of
non-Archimedean entire functions, \emph{Contemporary Mathematics}
\textbf{551} (2011), 57--69.
\url{https://arxiv.org/abs/1007.0984v2}.

\bibitem{vdH}
J. van der Hoeven, Operators on generalized power series,
\emph{Illinois Journal of Mathematics} \textbf{45} (2001), no.~4,
1161--1190. DOI:
\href{https://doi.org/10.1215/ijm/1258138061}{10.1215/ijm/1258138061}.
\url{https://www.texmacs.org/joris/noeth/noeth.html}.

\bibitem{ADH}
M. Aschenbrenner, L. van den Dries, and J. van der Hoeven,
On numbers, germs, and transseries, arXiv:1711.06936,
revised version. \url{https://arxiv.org/abs/1711.06936}.
Used for bibliographic cross-checks of the surreal background.

\bibitem{RepoRoot}
V. Reshetnikov, \emph{Surreal}: repository README and reported
formalization scope, inspected 21 September 2026.
\href{https://github.com/VladimirReshetnikov/Surreal/blob/4896a2808ce30e01b1c86ae3ba2295a64762d246/README.md}{Repository README, pinned snapshot}.

\bibitem{RepoCatalogue}
V. Reshetnikov, \emph{The surreal and surcomplex reports}, documentation
catalogue in the supplied \emph{Surreal} repository, inspected
21 September 2026. AI-assisted research collection.
\href{https://github.com/VladimirReshetnikov/Surreal/blob/4896a2808ce30e01b1c86ae3ba2295a64762d246/docs/README.md}{Documentation catalogue, pinned snapshot}.

\bibitem{RepoRankOne}
\emph{Rank-One Non-Archimedean Analytic Geometry inside the Surcomplex
Numbers}, research report in the supplied \emph{Surreal} repository;
README and scope inspected 21 September 2026. AI-assisted draft.
\href{https://github.com/VladimirReshetnikov/Surreal/tree/4896a2808ce30e01b1c86ae3ba2295a64762d246/docs/surcomplex/rank-one-berkovich}{Rank-one report package, pinned snapshot}.

\bibitem{RepoDivisors}
\emph{Global Support Obstructions: Moving Divisors, Mittag--Leffler,
and a Picard Dichotomy}, research report in the supplied \emph{Surreal}
repository; README and scope inspected 21 September 2026. AI-assisted
draft.
\href{https://github.com/VladimirReshetnikov/Surreal/tree/4896a2808ce30e01b1c86ae3ba2295a64762d246/docs/surcomplex/global-divisors}{Global-divisor report package, pinned snapshot}.
\end{thebibliography}
\end{document}
